\documentclass[11pt,a4paper]{article}

\usepackage[T1]{fontenc}
\usepackage{lmodern}
\usepackage{amsmath,amssymb,amsthm,mathtools}
\usepackage[margin=1in]{geometry}
\usepackage[hidelinks]{hyperref}

\newtheorem{theorem}{Theorem}[section]
\newtheorem{lemma}[theorem]{Lemma}
\newtheorem{proposition}[theorem]{Proposition}
\newtheorem{corollary}[theorem]{Corollary}
\theoremstyle{definition}
\newtheorem{definition}[theorem]{Definition}
\newtheorem{example}[theorem]{Example}
\theoremstyle{remark}
\newtheorem{remark}[theorem]{Remark}

\newcommand{\R}{\mathbb{R}}
\newcommand{\N}{\mathbb{N}}

\title{Mathematical Notes: Topic}
\author{Yibin Shen}
\date{\today}

\begin{document}
\maketitle

\section{Introduction}

Write a short summary of the topic and its references.

\section{Statements and proofs}

\begin{definition}
A function $f:\R\to\R$ is \emph{nonnegative} if
$f(x)\geq 0$ for every $x\in\R$.
\end{definition}

\begin{proposition}
If $f,g:\R\to\R$ are nonnegative, then $f+g$ is nonnegative.
\end{proposition}

\begin{proof}
For each $x\in\R$, we have $f(x)\geq 0$ and $g(x)\geq 0$.
Hence $(f+g)(x)=f(x)+g(x)\geq 0$.
\end{proof}

\section*{References}
% Replace with full bibliographic references.

\end{document}
